\section{Verlinde 公式：从共形块到融合规则}
\label{sec:verlinde}

用户原话：“这个公式怎么从共形块推出来？它的物理意义是什么？”

这一节先写公式是什么，再写 Witten 的推导链，最后写它的地位——因为这条公式是本文里“物理推导”与“数学定理”距离最近、也最容易被混淆的地方。

\subsection{公式}

对可积表示 $i,j,k$（$SU(2)$ 时 $2j\le k$），融合系数
\begin{equation}
N_{ijk}=\sum_l \frac{S_{il}S_{jl}S_{kl}}{S_{0l}} ,
\label{eq:verlinde}
\end{equation}
其中 $S$ 是模变换矩阵，由 $q=\exp\frac{2\pi\ii}{k+N}$ 决定 \src{382}。$N_{ijk}$ 有三重身份：
\begin{enumerate}[leftmargin=2em]
\item 融合规则 $R_i\otimes R_j=\bigoplus_k N_{ij}{}^k R_k$ 里的重数；
\item $S^2$ 上三个标记点（表示 $R_i,R_j,R_k$）的 Hilbert 空间维数，$\dim\cH_{S^2}(P_i,P_j,P_k)=N_{ijk}$ \src{373--374}；
\item 三点函数空间的维数 $\dim\Inv(R_i\otimes R_j\otimes R_k)$。
\end{enumerate}

\subsection{推导链：十个步骤，三处输入}

Witten 在 §4.5 “Some Concrete Surgeries”（起 p.388）里用一段斜体小标题 “Proof of Verlinde's Conjecture” 展开推导 \src{391}。注意它是\textbf{段首小标题}，不是节标题——这也提示读者：这一段是本文对 Verlinde 猜想的“物理重证”，而猜想本身已由 Moore--Seiberg 证明 \src{391}。推导链可以压缩成十步，但\textbf{真正的输入只有三处} \ours：

\begin{enumerate}[leftmargin=2em]
\item \textbf{手术律} (4.27) \src{385}／(4.28) \src{386}：沿链环 $C$ 做手术时，配分函数按 $S$ 矩阵展开。这是链的动力学输入。
\item \textbf{$N_{abc}$ 的定义} (4.36) \src{388}：把 $N_{abc}$ 定义为 $S^2\times S^1$ 上三标记点的配分函数，也就是三维 Hilbert 空间的维数。这是“把要求解的量翻译成可算的量”。
\item \textbf{秩一性} (4.49) \src{393}：$\dim\cH=1$ 时的归一化。这是最后一步把比值定下来的条件。
\end{enumerate}

中间的关键一步是第 6 步：把“平行线”变成“链接线”，即把 $\delta_{ij}$ 换成 $S_{ij}$——这一步用的是广义手术公式 (4.28)，也是全链唯一需要它的地方 \ours。

结果是 \src{393}（(4.51)）：
\begin{equation}
\sum_k S_i{}^k\,N_{kl}{}^m=\frac{S_{il}S_i{}^m}{S_{0i}} .
\label{eq:verlinde-key}
\end{equation}

\begin{keypoint}
\textbf{这条推导链没有用到共形块模型的任何动力学}。它只用“配分函数在手术下如何变换”。所以 (4.51) 的内容完全可以从三维理论的自洽性推出——这正是 Verlinde 公式“物理意义”的核心：\textbf{融合规则不是额外输入，它被模变换矩阵 $S$ 完全决定}。
\end{keypoint}

\subsection{两种读法}

(4.51) 可以朝两个方向读 \ours：

\begin{itemize}[leftmargin=1.6em]
\item \textbf{读法 (a)：融合规则 $\to S$}。若已知 $N_{ij}{}^k$，则 (4.51) 说 $S$ 把它们对角化。这是“陈述”。
\item \textbf{读法 (b)：$S\to$ 幂等元}。定义 $w_i=S_{0i}\sum_m S_i{}^m v_m$，则 (4.51) 等价于
\begin{equation}
w_i w_j=\delta_{ij}w_j ,
\end{equation}
即 $w_i$ 是 Verlinde 代数 $v_i v_j=\sum_k N_{ij}{}^k v_k$ 的一组正交幂等元，$N_{ijk}$ 由 $S$ 显式给出。这是“构造”，也是 Moore--Seiberg 之后的通行读法 \src{393}。
\end{itemize}

两种读法数学等价，但解释不同：(a) 把 $S$ 当作派生量，(b) 把 $S$ 当作基本量、融合规则是它的影子。物理上 (b) 更自然——$S$ 来自模变换，是几何的；融合规则是代数后果。

\subsection{物理意义}

一句话：\textbf{在三维 Chern--Simons 理论里，二维融合规则被三维的模变换完全决定}。所以：

\begin{itemize}[leftmargin=1.6em]
\item Verlinde 公式不是“又一个公式”，而是“三维理论自洽性对二维数据的约束”；
\item 它给出了 $N_{ijk}$ 的算法：只要知道 $S$，不用解任何共形场论就能算融合系数；
\item 反过来，它把 $S$ 矩阵变成主角——这也解释了为什么后续的模张量范畴、Reshetikhin--Turaev 不变量都以 $S$ 为中心。
\end{itemize}

\subsection{地位：物理推导 vs 数学定理}

原文的说法是：Verlinde 提出猜想 [12]，Moore--Seiberg 给出证明 [13]，“我们给出新证明” \src{391}。这里的“证明”是物理推导——用配分函数在手术下的形式性质，不是从 CS 测度出发的数学证明。

数学侧的四级路线要分清 \src{Schottenloher 书 p.177；上野 p.106}：
\begin{center}
\small
\begin{tabular}{@{}p{3.4cm}p{4.2cm}p{5.4cm}@{}}
\toprule
层面 & 谁 & 状态与工具 \\
\midrule
物理推导 & Witten §4.5（本文） & TQFT 公理下的推导，自洽可复算 \\
物理证明 & Moore--Seiberg 1989 & 首证；Witten 自认“新证明” \\
数学证明（几何） & TUY 1989；Faltings 1994 & 共形块丛 + KZ/Hitchin；rank-level 对偶 \\
组合/范畴 & RT 1991；模张量范畴 & 独立路线，以 $S$ 为中心 \\
现代代数证明 & 融合环 = Frobenius 代数 & 谱分解，把“自然数”变成代数事实 \\
\bottomrule
\end{tabular}
\end{center}

Schottenloher 的原话是“这些证明都不容易理解”：KZ 单值性加非交换 Hodge 理论是重机器；右端是 $\sin$ 的幂次和，“是自然数”本身都不显然 \src{Schottenloher 书 p.177}。

\begin{keypoint}
\textbf{一个常见范畴错误}：河野书的命题 1.22（书 p.37）与定理 2.16（书 p.60）用的是同一个 $S$ 矩阵，但前者是融合代数的基变换（代数侧，就是 Verlinde 公式），后者是 Dehn 手术核（几何侧，定义 Witten 不变量）。定理 2.16 的证明拿 1.22 当砖，但把手术公式当成“Verlinde 的证明”是张冠李戴 \src{河野 p.60--61}。
\end{keypoint}

\subsection{一个可当场做的验证}

$SU(2)_k$ 的 $S$ 矩阵有显式公式
\begin{equation}
S_{jl}=\sqrt{\frac{2}{k+2}}\,\sin\frac{\pi(2j+1)(2l+1)}{k+2},
\qquad j,l=0,\tfrac12,\dots,\tfrac{k}{2}.
\label{eq:smatrix}
\end{equation}
它由 $q=\exp\frac{2\pi\ii}{k+2}$ 与 (5.29)(5.30) 给出 \src{上野 p.103；河野 p.60}。

\textbf{$k=1$}：只有 $j=0,\frac12$，
\begin{equation}
S=\frac{1}{\sqrt2}\begin{pmatrix}1&1\\1&-1\end{pmatrix}.
\end{equation}
代入 \eqref{eq:verlinde} 逐个算：
\begin{align}
N_{\frac12,\frac12,0}&=\sum_l\frac{S_{\frac12 l}S_{\frac12 l}S_{0l}}{S_{0l}}
=\sum_l S_{\frac12 l}^2=\frac12+\frac12=1,\\
N_{\frac12,\frac12,\frac12}&=\sum_l\frac{S_{\frac12 l}^3}{S_{0l}}
=\frac{(1/\sqrt2)^3}{1/\sqrt2}+\frac{(-1/\sqrt2)^3}{1/\sqrt2}
=\frac12-\frac12=0 .
\end{align}
后一式正是截断的直接后果：$\frac12\otimes\frac12=0$，而 $j=1$ 的耦合被 $2j\le k$ 切掉，所以三点耦合为零。\textbf{注意}：这里“$N_{111}=0$”里的 $1$ 指 $j=1$（不可积），不要与“$j=\frac12$ 记作 1”混用。

\textbf{$k=2$}：$j=0,\frac12,1$，由 \eqref{eq:smatrix} 得 $S_{00}=\frac12$、$S_{0\frac12}=\frac{1}{\sqrt2}$、$S_{01}=\frac12$、$S_{\frac12\frac12}=0$、$S_{\frac12 1}=-\frac{1}{\sqrt2}$、$S_{11}=\frac12$。于是
\begin{equation}
N_{\frac12,\frac12,0}=1,\qquad N_{\frac12,\frac12,1}=1,
\end{equation}
即 $\frac12\otimes\frac12=0\oplus1$；同法得 $\frac12\otimes1=\frac12$、$1\otimes1=0\oplus1$。与截断融合规则逐条对上 \ours（这组数由 S5-B 用 python 独立复算过）。

\textbf{这个例子说明两件事}：一是 Verlinde 公式确实“只要 $S$ 就能算融合”；二是可积性截断必须同时作用在求和指标与结果指标上，否则会得到非整数。


\begin{boundary}
\begin{itemize}[leftmargin=1.4em]
\item 原文给出：§4.5 的十步链、(4.51)、两种读法、(4.36) 的定义、$SU(2)$ 的 $q$；段首小标题 “Proof of Verlinde's Conjecture” 在 p.391，属 §4.5（节标题 Some Concrete Surgeries，p.388）。
\item 我们补的：三处输入的归纳；两种读法的方向区分；$SU(2)_{k=1,2}$ 的数值验证；河野 1.22 与 2.16 的范畴区分。
\item 未证：数学严格证明（TUY/Faltings 只按二手转述，未读原文）；$S$ 矩阵归一化的细节；标架反常 (4.41) 在一般手术下的追踪——本文主链只用标准 $S$，Witten 明说 (4.41) 那种修正“在本文的简单应用中不必要” \src{390--391}。
\end{itemize}
\end{boundary}
